\section{The Abel map}

\begin{definition}[Simple sheaf]
\label{akcp-5-1-simple}
\dcref{5.1}
\source{MR0555258-compactifying-picard:page-0033}
An $\OO_X$-module $\II$ is simple over $S$ if it is locally finitely presented,
$S$-flat, and for every $T\to S$ the canonical map
$\OO_T\to(f_T)_*\mathcal H\!om(\II_T,\II_T)$ is an isomorphism.
\end{definition}

\begin{proposition}[Simple locus]
\label{akcp-5-2-simple-locus}
\dcref{5.2}
\source{MR0555258-compactifying-picard:page-0033}
\uses{akcp-1-1-H,akcp-1-3-local-free,akcp-5-1-simple}
For proper finitely presented $f:X\to S$ and locally finitely presented
$S$-flat $\II$, there is an open retrocompact $U\subset S$ such that
$T\to S$ factors through $U$ iff $\II_T$ is $T$-simple.
\end{proposition}
\begin{proof}
Put $H=H(\II,\II)$ and let $u:H\to\OO_S$ correspond to the identity.
Fiberwise $u$ is nonzero and hence surjective; its kernel is finitely generated
and commutes with base change.  Set $U=S\setminus\Supp\ker u$.  On $U$, $u$
is an isomorphism, and the base-change property of $H$ gives the required
scalar endomorphism isomorphisms for every further $T$.  Conversely, if
$\II_T$ is simple, the fiber of $\ker u$ at every point of $T$ vanishes;
Nakayama and finite generation force $T$ to factor through $U$.  Retrocompactness
follows by noetherian approximation.
\end{proof}

\begin{corollary}[Fiber criterion for simplicity]
\label{akcp-5-3-simple-fibers}
\dcref{5.3}
\source{MR0555258-compactifying-picard:page-0034}
\uses{akcp-5-2-simple-locus}
Under the hypotheses of \cref{akcp-5-2-simple-locus}, $\II$ is $S$-simple iff
$\OO_{k(s)}\to\Hom_{X(s)}(\II(s),\II(s))$ is an isomorphism for every point
(equivalently every geometric point) $s$.
\end{corollary}
\begin{proof}
Each condition says exactly that the simple locus of
\cref{akcp-5-2-simple-locus} is all of $S$.
\end{proof}

\begin{lemma}[Simplicity on an integral proper scheme]
\label{akcp-5-4-integral-simple}
\dcref{5.4}
\source{MR0555258-compactifying-picard:page-0034}
If $X$ is proper and irreducible over an algebraically closed field, satisfies
$S_1$, and a coherent $\II$ is generically $\OO_{X,\eta}$, then $\II$ is simple.
\end{lemma}
\begin{proof}
The $S_1$ condition puts $\Hom(\II,\II)$ inside its generic fiber, which is
the field $K(X)$.  It is therefore a domain.  Properness makes it finite
dimensional over the algebraically closed ground field, so it equals that field.
\end{proof}

\begin{definition}[The simple-sheaf functor]
\label{akcp-5-5-spl}
\dcref{5.5}
\source{MR0555258-compactifying-picard:page-0034}
Let $\Spl_{X/S}(T)$ be isomorphism classes of $T$-simple sheaves on $X_T$,
where $\II\sim\JJ$ if $\II\otimes\NN\simeq\JJ$ for an invertible $\OO_T$-module
$\NN$.  Write $\Spl_{X/S}^{\tau}$ for the associated sheaf in topology
$\tau\in\{\mathrm{Zar},\acute{\mathrm et},\mathrm{fppf},\mathrm{fpqc}\}$.
\end{definition}

\begin{proposition}[Separatedness and geometric points]
\label{akcp-5-6-spl-sheaf}
\dcref{5.6}
\source{MR0555258-compactifying-picard:page-0035}
\uses{akcp-4-3-twist-descent,akcp-5-5-spl}
The presheaf $\Spl_{X/S}$ is separated for the fpqc topology.  Consequently
there are canonical monomorphisms
$\Spl^{\mathrm{Zar}}\hookrightarrow\Spl^{\acute{\mathrm et}}
\hookrightarrow\Spl^{\mathrm{fppf}}\hookrightarrow\Spl^{\mathrm{fpqc}}$.
For a geometric point $t$, every $k(t)$-point of the fppf sheaf is represented
by a simple sheaf on $X(t)$.
\end{proposition}
\begin{proof}
If two twists become isomorphic fpqc-locally, the implication
(d)$\Rightarrow$(b) of \cref{akcp-4-3-twist-descent} descends the twist, giving
separatedness.  Sheafification preserves monomorphisms.  For the geometric
point assertion, an fppf representative over an extension algebra has a
$k(t)$-rational point because $k(t)$ is algebraically closed; pull back there.
\end{proof}

\begin{definition}[Hilbert-polynomial components]
\label{akcp-5-7-spl-polynomial}
\dcref{5.7}
\source{MR0555258-compactifying-picard:page-0035}
For projective $X/S$ with relatively very ample $\OO_X(1)$ and polynomial
$\varphi$, let $\Spl^{\varphi}_{X/S}$ be the subfunctor of
$\Spl^{\acute{\mathrm et}}_{X/S}$ whose local representatives have fiber
Hilbert polynomial $\varphi$.
\end{definition}

\begin{lemma}[Polynomial decomposition]
\label{akcp-5-8-spl-polynomial-open}
\dcref{5.8}
\source{MR0555258-compactifying-picard:page-0035}
\uses{akcp-5-6-spl-sheaf,akcp-5-7-spl-polynomial}
Each $\Spl^{\varphi}_{X/S}$ is open and closed in
$\Spl^{\acute{\mathrm et}}_{X/S}$, and these subfunctors cover it.
\end{lemma}
\begin{proof}
After an etale cover carrying a representative $\II$, the locus where the
fiber Hilbert polynomial equals $\varphi$ is open and closed by cohomology and
base-change.  Its image under the faithfully flat etale cover is open and
closed and represents the fiber product with the corresponding component.
Every coherent fiber has some polynomial, giving the covering.
\end{proof}

\begin{definition}[Relative torsion-free and pseudo-invertible]
\label{akcp-5-9-relative-tf}
\dcref{5.9}
\source{MR0555258-compactifying-picard:page-0036}
When the geometric fibers of $X/S$ are integral, an $\OO_X$-module is
relatively torsion-free rank one (resp. relatively pseudo-invertible) if it is
locally finitely presented and $S$-flat and each geometric fiber is rank-one
torsion-free (resp. rank-one torsion-free and Cohen--Macaulay).
\end{definition}

\begin{proposition}[Relative sheaves are simple]
\label{akcp-5-10-relative-simple}
\dcref{5.10}
\source{MR0555258-compactifying-picard:page-0036}
\uses{akcp-5-3-simple-fibers,akcp-5-4-integral-simple,akcp-5-9-relative-tf}
For proper finitely presented $X/S$ with integral geometric fibers, every
relatively torsion-free rank-one or relatively pseudo-invertible module is
$S$-simple.
\end{proposition}
\begin{proof}
Apply the fiber criterion \cref{akcp-5-3-simple-fibers}; each fiber satisfies
the hypotheses of \cref{akcp-5-4-integral-simple}.
\end{proof}

\begin{definition}[Picard subfunctors]
\label{akcp-5-11-picard}
\dcref{5.11}
\source{MR0555258-compactifying-picard:page-0036}
Let $\Pic^{\mathrm{CM}}_{X/S}$ and $\Pic^{\mathrm{tf}}_{X/S}$ be the subfunctors
of $\Spl_{X/S}$ represented by relatively pseudo-invertible and relatively
torsion-free rank-one sheaves, respectively.  Intersecting with a polynomial
component gives $\Pic^{\varphi,\mathrm{CM}}$ and
$\Pic^{\varphi,\mathrm{tf}}$.
\end{definition}

\begin{lemma}[Openness of fiber conditions]
\label{akcp-5-12-open-fiber-conditions}
\dcref{5.12}
\source{MR0555258-compactifying-picard:page-0036}
\uses{akcp-3-1-rank-one,akcp-5-9-relative-tf}
For a proper finitely presented $f$ and flat finitely presented $\II$, the
locus where $\II(s)$ is invertible is open retrocompact.  If fibers are
geometrically integral of fixed dimension, the loci where $\II(s)$ is
rank-one torsion-free, respectively pseudo-invertible, are open retrocompact.
\end{lemma}
\begin{proof}
Invertibility and Cohen--Macaulayness are open by cohomology and depth
semicontinuity.  Geometric reducedness and pure support dimension are open;
combine these with the characterization in \cref{akcp-3-1-rank-one}.  The
standard noetherian approximation argument gives retrocompactness.
\end{proof}

\begin{proposition}[Picard loci are open]
\label{akcp-5-13-picard-open}
\dcref{5.13}
\source{MR0555258-compactifying-picard:page-0037}
\uses{akcp-5-9-relative-tf,akcp-5-12-open-fiber-conditions}
If $\OO_X$ is $S$-simple, $\Pic^{\mathrm{CM}}_{X/S}$ is an open retrocompact
subsheaf of $\Spl^{\acute{\mathrm et}}_{X/S}$.  If integral fibers have fixed
dimension, $\Pic^{\mathrm{CM}}$ is open retrocompact in $\Pic^{\mathrm{tf}}$ and
$\Pic^{\mathrm{tf}}$ is open retrocompact in $\Spl^{\acute{\mathrm et}}$.
\end{proposition}
\begin{proof}
Represent a point etale-locally by $\II$ and apply
\cref{akcp-5-12-open-fiber-conditions}; flat locally finitely presented maps
are open, and the images of the loci represent the indicated fiber products.
Retrocompactness is checked on quasi-compact etale pieces.
\end{proof}

\begin{definition}[Quotient strata]
\label{akcp-5-14-quot-strata}
\dcref{5.14}
\source{MR0555258-compactifying-picard:page-0038}
For a locally finitely presented $\FF$ on proper finitely presented $X/S$,
$\Spl(\FF)$, $Q\operatorname{-div}(\FF)$, and
$P\operatorname{-div}(\FF)$ are the subfunctors of $\Quot_{\FF/X/S}$ whose
pseudo-ideals are respectively simple, relatively torsion-free rank one, and
relatively pseudo-invertible.  Add a superscript $\varphi$ for a fixed Hilbert
polynomial component.
\end{definition}

\begin{proposition}[Representability of quotient strata]
\label{akcp-5-15-quot-strata-representable}
\dcref{5.15}
\source{MR0555258-compactifying-picard:page-0038}
\uses{akcp-2-7-local-quot,akcp-5-2-simple-locus,akcp-5-12-open-fiber-conditions,akcp-5-14-quot-strata}
For locally projective $X/S$ with integral geometric fibers, all three quotient
strata are represented by retrocompact open subschemes of the corresponding
Quot scheme.  If $X/S$ is strongly projective and $\FF$ is a quotient of
$f^*B\otimes\OO_X(n)$ as in \cref{akcp-2-6-quot-representability}, each fixed
polynomial stratum is strongly quasi-projective.
\end{proposition}
\begin{proof}
The defining fiber conditions are open by
\cref{akcp-5-2-simple-locus} and \cref{akcp-5-12-open-fiber-conditions}; intersect with
the representable Quot strata of \cref{akcp-2-7-local-quot}.  Strong
quasi-projectivity follows from the open retrocompact inclusion into the
strongly quasi-projective Quot scheme.
\end{proof}

\begin{definition}[Abel map]
\label{akcp-5-16-abel-map}
\dcref{5.16}
\source{MR0555258-compactifying-picard:page-0038}
The Abel map associated with $\FF$ sends a quotient $\FF_T\twoheadrightarrow\GG$
to the class of its pseudo-ideal $\II(\GG)$ in $\Spl_{X/S}$.  The fiber over a
 simple sheaf $\II$ is the linear-system functor
$\LinSyst(\II,\FF)$.
\end{definition}

\begin{lemma}[Fibers of the Abel map]
\label{akcp-5-17-abel-fiber}
\dcref{5.17}
\source{MR0555258-compactifying-picard:page-0039}
\uses{akcp-4-2-linear-representability,akcp-5-6-spl-sheaf,akcp-5-16-abel-map}
Let $\II$ be $T$-simple and $\FF$ $S$-flat.  The fiber of the Abel map over
$[\II]$ is an open retrocompact subscheme $U$ of
$\PP(H(\II,\FF_T))$, with universal exact sequence
$0\to\II\otimes\OO_U(-1)\to\FF_U\to\GG\to0$.
If fibers are integral and $\II,\FF$ relatively torsion-free rank one, then
$U=\PP(H(\II,\FF_T))$.
\end{lemma}
\begin{proof}
The Cartesian square follows from separatedness of the simple-sheaf sheaf and
the definition of the Abel map.  Apply \cref{akcp-4-2-linear-representability}
to obtain $U$ and the universal sequence.  On an integral fiber every nonzero
map between rank-one torsion-free modules is injective by
\cref{akcp-3-4-hom-rank-one}, so the final assertion follows from the criterion
in the linear-system theorem.
\end{proof}

\begin{theorem}[Fiber dimension and smoothness]
\label{akcp-5-18-abel-smooth}
\dcref{5.18}
\source{MR0555258-compactifying-picard:page-0039}
\uses{akcp-1-3-local-free,akcp-5-17-abel-fiber}
Let $t$ be a geometric point of the simple-sheaf sheaf, represented by $\II$.
If every nonzero $\II\to\FF(t)$ is injective, then
\[
\dim \alpha^{-1}(t)=\dim_{k(t)}\Hom(\II,\FF(t))-1.
\]
If $\Ext^1_{X(t)}(\II,\FF(t))=0$, the Abel map is smooth along this fiber.
\end{theorem}
\begin{proof}
The fiber is the open $U$ of the projective space in
\cref{akcp-5-17-abel-fiber}; it is empty precisely when the Hom space is zero,
and otherwise has the displayed dimension.  The vanishing hypothesis and
\cref{akcp-1-3-local-free} make $H(\II,\FF)$ locally free near $t$, so the
projective bundle is smooth.  Smoothness descends through the etale cover used
to represent $t$.
\end{proof}

\begin{lemma}[Uniform regularity for a torsion-free family]
\label{akcp-5-19-uniform-regularity}
\dcref{5.19}
\source{MR0555258-compactifying-picard:page-0040}
\uses{akcp-3-4-hom-rank-one}
For a projective family with integral fibers and a relatively rank-one
torsion-free $\FF$ of one Hilbert polynomial, all fibers are $m$-regular for a
uniform $m$, $f_*\FF(m)$ is locally free of rank $\varphi(m)$, and the canonical
map $f^*f_*\FF(m)\otimes\OO_X(-m)\to\FF$ is surjective.
\end{lemma}
\begin{proof}
Apply the uniform regularity assertion of
\cref{akcp-3-4-hom-rank-one} with $\OO_X$ as the first rank-one sheaf.  Global
generation and cohomology-and-base-change for an $m$-regular family give the
last two claims.
\end{proof}

\begin{theorem}[Properness and finiteness of the Abel map]
\label{akcp-5-20-abel-proper}
\dcref{5.20}
\source{MR0555258-compactifying-picard:page-0040}
\uses{akcp-2-6-quot-representability,akcp-5-15-quot-strata-representable,akcp-5-17-abel-fiber,akcp-5-19-uniform-regularity}
Let $\FF$ be relatively rank-one torsion-free and let $P$ be a subsheaf of
$\Pic^{\acute{\mathrm et}}_{X/S}$.  The restriction $\alpha^{-1}(P)\to P$ is
proper and finitely presented.  If $X/S$ and $\FF$ have fixed Hilbert
polynomials, each polynomial component is strongly projective.  If $P$ has a
universal sheaf, this restriction is the structure morphism of the
corresponding quotient stratum.
\end{theorem}
\begin{proof}
Etale-locally on $P$, the fiber is the projective-space open of
\cref{akcp-5-17-abel-fiber}; in the torsion-free rank-one case it is the whole
projective space.  Properness and finite presentation descend from this local
description.  For fixed Hilbert polynomials, use
\cref{akcp-5-19-uniform-regularity} and \cref{akcp-2-6-quot-representability}
to embed the source in a projective bundle; a proper finitely presented
subscheme is strongly projective.  A universal sheaf identifies the source with
the quotient stratum by the same fiber description.
\end{proof}

\begin{remark}[Universal sheaves]
\label{akcp-5-21-universal-sheaf}
\dcref{5.21}
\source{MR0555258-compactifying-picard:page-0041}
\uses{akcp-5-20-abel-proper}
The existence of a universal sheaf on an entire Picard component is stronger
than representability by an etale sheaf; it is equivalent to that component
already being an etale sheaf.  A universal sheaf exists after rigidification,
for example when the smooth locus admits a section.
\end{remark}
